\clearpage\thispagestyle{empty}
\pdfbookmark[0]{Quellen zu Komet (17. Mai 1821)}{cometmay1821-sources}
\begin{center}\textbf{Quellen zu Komet (17. Mai 1821)}\end{center}
\noindent Diese KI-Übertragung ins heutige Deutsch folgt dem Bericht \textit{Comet.} in \href{https://gdz.sub.uni-goettingen.de/id/PPN236007467}{Gauß’ \textit{Werken}, Bd.~VI} (Göttingen, 1874), S.~437--439. Der Erstdruck erschien unter Göttingen in den \textit{Göttingischen gelehrten Anzeigen}, Nr.~78, 17. Mai 1821, S.~\href{https://images.p.sub.uni-goettingen.de/iiif/image/gdz:PPN319721507_1821:00000773/full/1600,/0/default.jpg}{769}, \href{https://images.p.sub.uni-goettingen.de/iiif/image/gdz:PPN319721507_1821:00000774/full/1600,/0/default.jpg}{770} und \href{https://images.p.sub.uni-goettingen.de/iiif/image/gdz:PPN319721507_1821:00000775/full/1600,/0/default.jpg}{771}. Das digitalisierte Exemplar stammt aus der Niedersächsischen Staats- und Universitätsbibliothek Göttingen. Der Bericht enthält die Göttinger Beobachtungen und nennt von Staudt als Urheber der parabolischen Elemente und der Helligkeitsrechnung.
\par\smallskip\noindent Beide Drucke nennen den Göttinger Meridian in der Zeile für den Logarithmus des kleinsten Abstands. Die Anordnung und die Deklination vom 10. Februar folgen der Werkausgabe. In der ersten Beobachtung verwendet der Erstdruck ein hochgestelltes u und Striche, die Werkausgabe hochgestellte u, m und s. Der Erstdruck verwendet Dezimalkommas. Der \hyperlink{native-440-part-lower}{Bericht vom 10. Februar 1821} ist eine frühere Göttinger Mitteilung über denselben Kometen. KI-Übertragung und KI-Satz.
\clearpage
